% Compile with: lualatex manual.tex
\documentclass[11pt,a4paper]{article}

\usepackage{fontspec}
\setmainfont{Segoe UI}
\setmonofont{Consolas}

\usepackage{geometry}
\geometry{margin=2.5cm}

\usepackage{xcolor}
\definecolor{accent}{HTML}{1E88E5}
\definecolor{darkbg}{HTML}{2D2D2D}
\definecolor{green}{HTML}{43A047}
\definecolor{red}{HTML}{E53935}

\usepackage{tcolorbox}
\tcbuselibrary{listings,skins}

\newtcolorbox{tipbox}{
  colback=accent!8, colframe=accent!60,
  fonttitle=\bfseries, title=Tip,
  left=4pt, right=4pt, top=2pt, bottom=2pt
}

\newtcolorbox{warnbox}{
  colback=red!8, colframe=red!60,
  fonttitle=\bfseries, title=Important,
  left=4pt, right=4pt, top=2pt, bottom=2pt
}

\newtcolorbox{cmdbox}{
  colback=darkbg, colframe=darkbg,
  colupper=white, fontupper=\ttfamily\small,
  left=6pt, right=6pt, top=4pt, bottom=4pt,
  sharp corners
}

\usepackage{enumitem}
\setlist[enumerate]{itemsep=3pt}
\setlist[itemize]{itemsep=3pt}

\usepackage{booktabs}
\usepackage{tabularx}
\usepackage{hyperref}
\hypersetup{colorlinks=true, linkcolor=accent, urlcolor=accent}

\usepackage{titlesec}
\titleformat{\section}{\Large\bfseries\color{accent}}{}{0pt}{}[\vspace{2pt}\hrule\vspace{4pt}]
\titleformat{\subsection}{\large\bfseries}{}{0pt}{}
\titleformat{\subsubsection}{\normalsize\bfseries}{}{0pt}{}

\pagestyle{plain}

\begin{document}

\begin{center}
{\fontsize{28}{34}\selectfont\bfseries\color{accent} Streamer}\\[6pt]
{\large Qobuz Music Downloader --- User Manual}\\[4pt]
{\small Version 0.2 --- June 2026}
\end{center}

\vspace{1em}
\tableofcontents

%=============================================================================
\section{Overview}
%=============================================================================

Streamer downloads music from Qobuz --- from 320\,kbps MP3 up to
24-bit/192\,kHz Hi-Res FLAC.  It comes in two forms that share the same
config file:

\begin{itemize}
  \item \textbf{streamer-gui.exe} --- graphical app for Windows.  Search, pick
        what you want, and download.  Works on any Windows~10+ machine.
  \item \textbf{streamer.exe} --- command-line tool.  Faster, scriptable, and
        runs anywhere you can open a terminal.
\end{itemize}

A valid Qobuz subscription is required.

%=============================================================================
\section{Login --- Getting Your Credentials}
%=============================================================================

Qobuz blocks direct email/password login from third-party tools.  Streamer
instead reads your session token from a browser where you are already logged
in, using a small browser script.  You only need to do this once (or whenever
your token expires, which is rare).

\subsection{Step 1 --- Install the Userscript}

\begin{enumerate}
  \item Install \textbf{Violentmonkey} (Firefox, Chrome, Edge) or
        \textbf{Tampermonkey} from your browser's extension store.
  \item Drag the file \texttt{streamer-login.user.js} (located in the same
        folder as \texttt{streamer-gui.exe}) onto the Violentmonkey dashboard
        tab, or open it directly in your browser.
  \item Click \textbf{Install} when the extension asks.
\end{enumerate}

\subsection{Step 2 --- Copy the Login Command}

\begin{enumerate}
  \item Open \texttt{https://play.qobuz.com} and log in normally --- with your
        email and password, Google, or Apple.  The script runs on that page.
  \item A button appears in the bottom-right corner of the page. It starts
        dimmed while waiting for your session to load:

        \medskip
        {\ttfamily\small\color{accent} [waiting] Waiting for login\ldots}
        \medskip

  \item Once the page has loaded your session (a second or two), the button
        becomes clickable:

        \medskip
        {\ttfamily\small\color{accent} [copy] Copy streamer login}
        \medskip

  \item Click it.  The script intercepts the authentication token that Qobuz
        sends with every page request and copies this command to your clipboard:
\end{enumerate}

\begin{cmdbox}
streamer.exe login --user-id "12345678" --token "eyJhb..."
\end{cmdbox}

\subsection{Step 3A --- Save via the GUI}

\begin{enumerate}
  \item Open \texttt{streamer-gui.exe} and go to the \textbf{Settings} tab.
  \item Click \textbf{+ Add} and select the new empty row.
  \item From the clipboard command, copy \emph{only the number} after
        \texttt{-{}-user-id} and paste it into the \textbf{User ID} field.
  \item Copy the long string after \texttt{-{}-token} and paste it into the
        \textbf{Token} field.
  \item Click \textbf{Save Credentials}.
\end{enumerate}

On success the row shows \textbf{Authenticated} and your country code (US, FR,
etc.)\ is detected automatically.

\begin{tipbox}
If you have multiple accounts (e.g.\ one US, one FR), repeat Steps~2--3 for
each account and add a separate row for each.
\end{tipbox}

\subsection{Step 3B --- Save via the CLI}

Paste the copied command directly into a terminal:

\begin{cmdbox}
streamer.exe login --user-id "12345678" --token "eyJhb..."
\end{cmdbox}

The tool contacts Qobuz to verify the token, detects your country code, and
writes everything to the config file.  You should see:

\begin{cmdbox}
Country: US\\
Login successful. Config saved.
\end{cmdbox}

\subsection{Step 4 --- Set Your Download Folder}

In the \textbf{Settings} tab, click \textbf{Browse\ldots} next to
\emph{Download Dir}, pick a folder, and click \textbf{Save Settings}.

Or from the terminal:

\begin{cmdbox}
streamer.exe config set-dir "D:\textbackslash{}Music"
\end{cmdbox}

\subsection{Step 5 --- Choose a Default Quality}

\begin{center}
\begin{tabular}{@{}llr@{}}
\toprule
\textbf{Setting}  & \textbf{Format}            & \textbf{Approx.\ size / album} \\
\midrule
\texttt{mp3}        & 320\,kbps MP3              & $\sim$100\,MB \\
\texttt{flac}       & FLAC 16-bit / 44.1\,kHz   & $\sim$300\,MB \\
\texttt{flac-hi}    & FLAC 24-bit / 96\,kHz      & $\sim$1\,GB   \\
\texttt{flac-ultra} & FLAC 24-bit / 192\,kHz     & $\sim$2\,GB   \\
\bottomrule
\end{tabular}
\end{center}

Scroll the \textbf{Quality} wheel in Settings and click \textbf{Save
Settings}.  You can override the quality per download from the CLI with
\texttt{-{}-quality}.

%=============================================================================
\section{Searching --- GUI}
%=============================================================================

Click the \textbf{Search} tab.

\subsection{How a Search Works}

\begin{enumerate}
  \item Type your query in the search box and press \textbf{Enter} (or click
        \textbf{Search}).  Streamer fetches up to 50 results per content type
        per account from Qobuz.
  \item Results appear in the table immediately.  A counter above the list
        shows \textbf{N\,/\,M results}: N is how many pass the current filter,
        M is the total fetched.
  \item \textbf{As you keep typing}, the visible results are filtered
        in real time --- without hitting the network again.  This lets you
        quickly narrow a large result set without a new search.
\end{enumerate}

\subsection{The Type Picker}

The scroll wheel on the right side of the search bar selects what kind of
content to fetch:

\begin{center}
\begin{tabular}{@{}ll@{}}
\toprule
\textbf{Mode}      & \textbf{What it does} \\
\midrule
Smart (auto)   & Detects the type from your query (see below) \\
Albums         & Only album results \\
Tracks         & Only track results \\
Artists        & Only artist results \\
Playlists      & Only playlist results \\
All types      & Fetch all four types at once \\
\bottomrule
\end{tabular}
\end{center}

\textbf{Smart (auto)} looks for a \texttt{type:} filter in your query and
uses that as the API type.  If no type filter is found, it fetches all types.
Example: typing \texttt{pink floyd type:album} on Smart mode only fetches
albums from the API --- faster than fetching everything.

\subsection{Filter Syntax}

Inside the search box you can combine free-text terms with structured filters.
Filters have the form \texttt{field:value} or, for numeric fields,
\texttt{field:>value}.

\subsubsection{String filters (partial match, case-insensitive)}

\begin{center}
\begin{tabularx}{\linewidth}{@{}l X l@{}}
\toprule
\textbf{Filter} & \textbf{Matches} & \textbf{Example} \\
\midrule
\texttt{artist:value}   & Artist name contains value  & \texttt{artist:"pink floyd"} \\
\texttt{title:value}    & Title contains value        & \texttt{title:"wish you were"} \\
\texttt{genre:value}    & Genre contains value        & \texttt{genre:jazz} \\
\texttt{label:value}    & Record label contains value & \texttt{label:ecm} \\
\texttt{type:value}     & Content type                & \texttt{type:album} \\
\texttt{country:value}  & Account country             & \texttt{country:US} \\
\bottomrule
\end{tabularx}
\end{center}

\subsubsection{Numeric filters (year and duration)}

\texttt{year} is the release year.  \texttt{duration} is in \textbf{seconds}.

\begin{center}
\begin{tabular}{@{}ll@{}}
\toprule
\textbf{Operator} & \textbf{Meaning} \\
\midrule
\texttt{year:2024}        & Exact year \\
\texttt{year:2020-2024}   & Year range (inclusive) \\
\texttt{year:>2020}       & After 2020 \\
\texttt{year:<=1975}      & Up to and including 1975 \\
\texttt{duration:>=600}   & At least 10 minutes (600 s) \\
\texttt{duration:<180}    & Under 3 minutes \\
\texttt{duration:300-600} & Between 5 and 10 minutes \\
\bottomrule
\end{tabular}
\end{center}

\subsubsection{Boolean filters}

\begin{center}
\begin{tabular}{@{}ll@{}}
\toprule
\textbf{Filter}         & \textbf{Meaning} \\
\midrule
\texttt{hires:true}     & Hi-Res audio available \\
\texttt{hires:false}    & Standard resolution only \\
\texttt{explicit:true}  & Explicit content \\
\texttt{explicit:false} & Clean / censored version \\
\bottomrule
\end{tabular}
\end{center}

\subsubsection{Logical operators}

Terms next to each other are AND by default.  You can also use:

\begin{center}
\begin{tabular}{@{}lll@{}}
\toprule
\textbf{Operator} & \textbf{Aliases} & \textbf{Example} \\
\midrule
AND (implicit) & \texttt{AND}, \texttt{y}, \texttt{\&\&} & \texttt{pink floyd hires:true} \\
OR             & \texttt{OR}, \texttt{o}, \texttt{||} & \texttt{artist:coltrane OR artist:mingus} \\
NOT            & \texttt{NOT}, \texttt{!}            & \texttt{jazz NOT explicit:true} \\
Parentheses    & \texttt{(} \texttt{)}               & \texttt{(artist:coltrane OR artist:mingus) type:album} \\
\bottomrule
\end{tabular}
\end{center}

\subsubsection{Spanish aliases}

All field names are also accepted in Spanish:
\texttt{artista}, \texttt{título}, \texttt{álbum}, \texttt{género},
\texttt{sello}, \texttt{tipo}, \texttt{país}, \texttt{año},
\texttt{duración}, \texttt{alta\_res}, \texttt{explícito}.

Boolean values are also accepted in Spanish: \texttt{sí} / \texttt{verdadero}
for \texttt{true}, and \texttt{no} / \texttt{falso} for \texttt{false}.

\subsection{Practical Examples}

\begin{cmdbox}
pancrasia type:album duration:>=600
\end{cmdbox}
Albums matching ``pancrasia'' with total duration at least 10 minutes.

\begin{cmdbox}
artist:"pink floyd" year:1970-1980 hires:true
\end{cmdbox}
Hi-Res Pink Floyd albums from the 1970s.

\begin{cmdbox}
(artist:"coltrane" OR artist:"mingus") type:track duration:>480
\end{cmdbox}
Tracks over 8 minutes by Coltrane or Mingus.

\begin{cmdbox}
género:jazz año:>2020 NOT explícito:sí
\end{cmdbox}
Clean jazz releases from after 2020.

\subsection{Reading the Results}

\begin{center}
\begin{tabular}{@{}ll@{}}
\toprule
\textbf{Column}  & \textbf{Content} \\
\midrule
Title    & Name, with version in parentheses when available \\
Artist   & Primary artist or performer \\
Label    & Record label \\
Date     & Full release date \\
Duration & Track or album total length (m:ss) \\
Genre    & Genre tag \\
Hi-Res   & \emph{Yes} if hi-res audio is available \\
E        & Explicit content marker \\
Type     & album / track / artist / playlist \\
Country  & Which of your accounts returned this result \\
\bottomrule
\end{tabular}
\end{center}

Results are grouped by artist, then sorted newest-first within each artist.
When you click a row, a short summary appears in the counter area above
the list.

%=============================================================================
\section{Downloading --- GUI}
%=============================================================================

\subsection{Downloading One Item}

Right-click any row and choose \textbf{Download}, or click a row to select it
and then click \textbf{Download Selected} at the bottom right.

\subsection{Downloading Multiple Items (the Cart)}

\begin{enumerate}
  \item \textbf{Tick the checkbox} on the left of each row you want.  Checked
        items appear in a \textbf{Selected} list below the results table and
        stay checked even if you run a new search.
  \item The button at the bottom right updates to show how many items are
        queued: \textbf{Download (3)}, for example.
  \item Click \textbf{Download (N)} to send all queued items to the
        \textbf{Downloads} tab.  The cart clears automatically.
  \item Click \textbf{Clear All} in the Selected section to uncheck everything
        without downloading.
\end{enumerate}

\begin{tipbox}
The cart survives new searches.  You can search for ``Coltrane'', tick some
albums, then search for ``Miles Davis'', tick more albums, and download all of
them at once.
\end{tipbox}

\subsection{What Gets Saved to Disk}

\begin{cmdbox}
Music/\\
  US/\\
~~~~Artist Name/\\
~~~~~~~~Album Title/\\
~~~~~~~~~~~~01. Track Title.flac\\
~~~~~~~~~~~~02. Track Title.flac\\
~~~~~~~~~~~~...\\
~~~~~~~~~~~~cover.jpg\\
~~~~~~~~~~~~album\_description.txt~~(if available)\\
~~~~~~~~~~~~booklet.pdf~~~~~~~~~~~~~(if available)\\
~~~~artist\_bio.txt~~~~~~~~~~~~~~~~~~(if available)\\
~~~~artist.jpg~~~~~~~~~~~~~~~~~~~~~~(if available)
\end{cmdbox}

If you have only one account with no country set, the country subfolder
(\texttt{US/}) is omitted.

\subsection{The Downloads Tab}

Switch here to watch progress.  The top part lists every queued and active
download with its current status.  The log at the bottom shows detailed
messages --- check it if something fails.

%=============================================================================
\section{Searching and Downloading --- CLI}
%=============================================================================

\subsection{Search}

\begin{cmdbox}
streamer.exe search "pink floyd" --type albums
\end{cmdbox}

Available types: \texttt{albums}, \texttt{tracks}, \texttt{artists},
\texttt{playlists}, \texttt{all}.  Default is \texttt{albums}.

\medskip
Limit the number of results (default 20):

\begin{cmdbox}
streamer.exe search "coltrane" --type tracks -n 50
\end{cmdbox}

\subsection{Duration Filter}

The CLI supports \texttt{-{}-min-duration} and \texttt{-{}-max-duration}.
When either flag is active, the duration is shown in brackets next to each
result (\texttt{[1:08:58]}).  Accepted formats:

\begin{center}
\begin{tabular}{@{}ll@{}}
\toprule
\textbf{Format}  & \textbf{Meaning} \\
\midrule
\texttt{600}     & 600 seconds (bare integer) \\
\texttt{90s}     & 90 seconds \\
\texttt{10m}     & 10 minutes \\
\texttt{2h}      & 2 hours \\
\texttt{1h30m}   & 1 hour 30 minutes \\
\texttt{1:30}    & 1 min 30 s (mm:ss) \\
\texttt{1:30:00} & 1 h 30 min (hh:mm:ss) \\
\bottomrule
\end{tabular}
\end{center}

\begin{cmdbox}
streamer.exe search "pancrasia" --type albums --min-duration 10m
\end{cmdbox}

\begin{cmdbox}
streamer.exe search "coltrane" --type tracks --min-duration 8m --max-duration 20m
\end{cmdbox}

\begin{warnbox}
Duration filtering happens after the API responds.  Streamer fetches up to
$5\times$ the requested limit to compensate, but if very few results in the
catalog match your range, you may receive fewer rows than \texttt{-n} asks for.
Raise \texttt{-n} if needed.
\end{warnbox}

Add \texttt{-{}-tsv} for tab-separated output (columns: id, title, artist,
label, date, \textbf{duration in seconds}, genre, hires, type, explicit):

\begin{cmdbox}
streamer.exe search "pancrasia" --type albums --min-duration 10m --tsv
\end{cmdbox}

\subsection{Download}

By URL:

\begin{cmdbox}
streamer.exe download "https://www.qobuz.com/us-en/album/..." --quality flac-hi
\end{cmdbox}

By bare ID (album IDs are alphanumeric strings; track IDs are integers):

\begin{cmdbox}
streamer.exe download abc123def\\
streamer.exe download 12345678 --quality flac
\end{cmdbox}

Use \texttt{-{}-account N} (0-based) to pick which account to download with.
The default is account 0.

%=============================================================================
\section{Multiple Accounts}
%=============================================================================

You can add as many Qobuz accounts as you like --- one per country is common
when catalogs differ by region.

\begin{itemize}
  \item In the GUI: click \textbf{+ Add} in Settings for each account and log
        in to each one separately.
  \item With the CLI: run \texttt{login} once per account.  Each call creates
        a new account slot (matched by user ID, so running it again for the
        same account just refreshes the token).
\end{itemize}

When you search in the GUI, Streamer queries \emph{every} configured account
and merges the results.  The \textbf{Country} column tells you which account
found each row, and the correct account is used automatically when you download.

Use \textbf{Export Accounts} / \textbf{Import Accounts} in Settings to back
up your credentials as a JSON file.

%=============================================================================
\section{Configuration File}
%=============================================================================

Location: \texttt{\%APPDATA\%\textbackslash{}streamer\textbackslash{}config.toml}

\begin{cmdbox}
{[[}accounts{]]}\\
country = "US"~~~~~~~~~~\# set automatically on login\\
email = ""~~~~~~~~~~~~~~\# optional, display only\\
app\_id = "..."\\
app\_secret = "..."\\
user\_id = "12345678"\\
auth\_token = "..."\\[6pt]
{[}settings{]}\\
download\_dir = "C:/Users/You/Music"\\
quality = "flac-ultra"~~\# mp3 | flac | flac-hi | flac-ultra\\
concurrency = 8~~~~~~~~~\# parallel download workers\\
requests\_per\_minute = 0 \# 0 = no limit
\end{cmdbox}

All fields are managed by the GUI and CLI.  You can edit this file by hand if
needed, but be careful to keep valid TOML syntax.

%=============================================================================
\section{Troubleshooting}
%=============================================================================

\begin{description}[style=nextline]

  \item[The script button stays dimmed after the page loads]
    The script captures tokens from network requests that happen \emph{after}
    the page starts.  If the tab was already open when you installed the
    script, reload the page (\textbf{F5}).  Make sure you are signed in to
    \texttt{play.qobuz.com} before clicking.

  \item[``Login successful'' but the account shows Not logged in in the GUI]
    The GUI reads the config file on startup.  Restart \texttt{streamer-gui.exe}
    after logging in from the CLI, or use the GUI's own credential fields.

  \item[Download fails with ``401 Unauthorized'' or ``Not authenticated'']
    Your session token has expired.  Repeat the login steps (Section~2) to
    get a fresh token.

  \item[Download fails with a path error (``os error 123'' on Windows)]
    The download directory path is invalid or does not exist.  Go to Settings,
    click Browse, pick an existing folder, and Save Settings.

  \item[Search returns nothing, or ``No account at index 0'']
    No account is configured.  Check the Settings tab: at least one row must
    show \textbf{Authenticated}.  If all show ``Not logged in'', follow
    Section~2 to log in again.

  \item[Duration filter returns fewer results than requested]
    Expected --- see the warning box in Section~5.2.  Increase \texttt{-n}.

  \item[File names contain garbled characters]
    This is a build issue: ensure \texttt{streamer-gui.exe} was compiled with
    the \texttt{/utf-8} flag (it is set in \texttt{CMakeLists.txt} by default).
    Re-download or rebuild.

\end{description}

\vfill
\begin{center}
\small\color{gray}
Streamer is not affiliated with Qobuz.\\
A valid Qobuz subscription is required to download music.
\end{center}

\end{document}
